\section{Background}
\subsection{Large Language Models}
\label{sect:background:preliminaries}

Large Language Models (LLMs) are a family of causal transformer models with multiple identically-structured layers. Each layer combines an attention block, a feed-forward network (FFN) and normalization layers. The input of each layer, $\mathbf{x}$, is an $N\times HD$ tensor, where $N$ is the number of input tokens, $H$ represents the number of attention heads, and $D$ is the hidden dimension for each head. Serving LLMs involves two stages: the \contextstage, where all prompt tokens are presented simultaneously ($N > 1$ for each request), and the \decodingstage, where the model only processes one token at a time for each prompt ($N = 1$ for each request).

In attention blocks, $\mathbf{x}$ first undergoes linear projection to obtain $\mathbf{q}\in\mathbb{R}^{N\times HD},\mathbf{k},\mathbf{v}\in\mathbb{R}^{N\times H_{KV}D}$, where $H_{KV}$ is the number of key/value heads. We have $H = H_{KV}$ in the standard multi-head attention (MHA), while recent methods~\cite{touvron2023llama2,jiang2023mistral,jiang2024mixtral} also employ grouped-query attention (GQA)~\cite{ainslie2023gqa} with $H = rH_{KV} (r\in \mathbb{Z})$. We concatenate $\mathbf{k},\mathbf{v}$ with pre-computed \textit{KV cache} features of $S$ previous tokens to obtain $\mathbf{K}, \mathbf{V}\in\mathbb{R}^{(S+N)\times H_{KV}D}$ and compute attention using:

\vspace{-10pt}
\begin{equation}
\small
\mathbf{o}_{h} = \text{softmax}\left(\frac{\mathbf{q}_{h}\mathbf{K}_{h_{KV}}^T}{\sqrt{D}}\right)\mathbf{V}_{h_{KV}},\quad h_{KV}=\left\lfloor\frac{h}{r}\right\rfloor
\label{eqn:background:attn}
\end{equation}
\vspace{-10pt}

The result $\mathbf{o}$ is multiplied with an output projection matrix $\mathbf{W}_O \in \mathbb{R}^{HD \times HD}$, and the product is added to $\mathbf{x}$ as the input of FFN. The FFN is composed of linear projection and activation layers and it does not mix features between tokens. 

\subsection{Integer Quantization}

Integer quantization maps high-precision numbers to discrete levels. The process can be formulated as: 


\begin{equation}
\resizebox{0.9\linewidth}{!}{
$\displaystyle 
\small
    \mathbf{Q}_{\mathbf{X}} = \left\lceil \frac{\mathbf{X}}{s}+z\right\rfloor, s=\frac{\mathbf{X}_{\max}-\mathbf{X}_{\min}}{q_{\max}-q_{\min}}, z= \left\lceil q_{\min}-\frac{\mathbf{X}_{\min}}{s}\right\rfloor
$%
}
\label{eqn:background:quantization}
\end{equation}
\vspace{-15pt}

where $\mathbf{X}$ is the floating point tensor, $\mathbf{Q}_{\mathbf{X}}$ is its $n$-bit quantized counterpart, $s$ is the scaling factor and $z$ is the zero point. Thus, the dequantized tensor can be represented as,

\vspace{-10pt}
\begin{equation}
\small
\hat{\mathbf{X}} = Q\left(\mathbf{X}\right) = \left(\mathbf{Q}_{\mathbf{X}} - z\right) \cdot s
\label{eqn:background:dequantization}
\end{equation}
\vspace{-18pt}

This is known as  \textit{asymmetric} quantization, where $\mathbf{X}_{\max} = \max\left(\mathbf{X}\right), \mathbf{X}_{\min} = \min\left(\mathbf{X}\right)$, and $q_{\max} - q_{\min} = 2^{n}-1$ for integer quantization. Equation \ref{eqn:background:quantization} can be further simplied to \textit{symmetric} quantization, where $z = \mathbf{0}$, $\mathbf{X}_{\max} = -\mathbf{X}_{\min} = \max\left|\mathbf{X}\right| $, and $q_{\max} - q_{\min} = 2^{n} -2$ .

In this paper, we denote $x$-bit weight, $y$-bit activation and $z$-bit KV cache quantization in LLMs as \textbf{\texttt{WxAyKVz}}, and use the abbreviated notation \textbf{\texttt{WxAy}} if \texttt{y=z}.  Apart from bit precision, quantization can also be applied at various granularities. \textit{Per-tensor} quantization shares $s$ and $z$ across the entire tensor. \textit{Per-channel} quantization for weights or \textit{per-token} quantization for activations means that $s$ and $z$ are shared within each \textit{row} of tensor. \textit{Per-group} quantization further reduces the degree of parameter sharing by using different $s$ and $z$ for every $g$ columns within each row, where $g$ is the group size.











